\section{Introduction}\label{introduction}

Autonomous ML experimentation agents can now propose code edits, run training jobs, inspect metrics, and choose follow-up trials. Evaluating such systems by final task score alone is insufficient. A higher validation AUC does not tell us whether the agent edited only allowed files, whether failed trials were preserved, whether no-op patches were rejected, whether keep/discard decisions were made by an external verifier, or whether the system repeatedly attempted known-bad ideas. For scientific ML, these process properties are not implementation details; they determine whether an autonomous experiment is auditable.

This paper argues that autonomous ML experimentation should be evaluated as a governed process. In the Agent = Model + Harness decomposition, the model proposes actions, but the harness supplies tools, memory, execution boundaries, state tracking, and verification. We instantiate that idea as a control plane for ML experimentation: a manager proposes one bounded change, a worker materialises the patch and run artifacts, and the control plane owns lifecycle, scope validation, metric parsing, keep/discard decisions, memory updates, and append-only records.

The central contribution is a governance-metric evaluation protocol for autonomous experimentation. We report acceptance rate, invalid-action rate, repeated-bad rate, provenance completeness, artifact capture completeness, and a failure taxonomy alongside the task metric. These metrics expose behaviours that task score hides: invalid edit attempts, missing metrics, repeated degraded proposals, no-op patches, stale pending trials, and unauditable decisions. A campaign with modest task improvement but complete lifecycle records is more scientifically interpretable than a campaign with an unexplained score increase. This framing targets the KDD Agentic AI Evaluation workshop's focus on measuring and benchmarking agent behaviour in real task settings.

We implement the protocol in \passthrough{\lstinline!autoresearch\_harness!}, a local harness around a controlled ResNet-trigger waveform-classification node. The node specification declares the editable surface (\passthrough{\lstinline!train.py!}), frozen data and pipeline files, the training command, metric parser, acceptance rule, and validity checks. The control plane then enforces a fixed-budget trial lifecycle: proposal, pending guard, patch/run, scope and no-op validation, metric parsing, deterministic decision, append-only ledger write, memory update, and guard release.

Empirically, we use the harness to run a five-trial real-worker campaign and optional strengthening runs on the same node. The primary campaign records two kept valid edits and three failed-invalid trials with complete provenance. Longer 10-trial memory runs show full lifecycle coverage and a clearer memory signal: the no-memory arm has repeated-bad rate 0.60, while both memory arms reduce it to 0.40 and reach a higher best validation AUC. A five-seed baseline replicate shows substantial task-metric variation, with mean validation AUC 0.784755 and bootstrap 95\% CI {[}0.774809, 0.798329{]}, so we treat AUC as secondary evidence rather than the headline claim.

This framing leads to three contributions:

\begin{itemize}
\tightlist
\item
  governance metrics as first-class evaluation criteria for autonomous ML experimentation, including acceptance, invalid-action, repeated-bad, provenance, artifact-completeness, and failure-taxonomy measures;
\item
  a governed control-plane protocol that separates proposal authority from execution, validation, and keep/discard authority;
\item
  an append-only audit schema and artifact pipeline that records proposal, patch, run log, parsed metric, decision, reproducibility hashes, and failure category for every trial, supported by stress failures, memory ablation, and seed/noise guardrails.
\end{itemize}

The rest of the paper is organised as follows. Section 2 positions the work relative to autonomous code-space exploration, scientific reproduction systems, harness engineering, experiment tracking, and AutoML. Section 3 describes the governed control plane. Section 4 defines the evaluation protocol and governance metrics. Section 5 reports campaign, ablation, stress, and reproducibility results. Section 6 discusses limitations and the path to broader multi-node evaluation.
